\scanpage{103}
\mspage{\unsure{90}}
% archive sheet number: 102

\deleted{\uline{\unsure{Définition}}} \inserted{\uline{Définition} 5.6.} Les triples de Quillen \unsure{déduits} \inserted{comme dans 5.5} d'une
catégorie à cofibrations spéciale sont dits
\uline{spéciaux}. Ainsi un \unsure{triples|triple} de Quillen
$(W,C,F)$ de $M$ est spécial s'il
satisfait les conditions a) à d) de la
définition (4.4.), concernant $(W,C)$ exclusivement, plus les deux conditions suivantes

\textcircled{\scriptsize $c'$} $\exists$ petite partie $TC_{0}$ de $TC$, telle
que $TC\subset (TC_{0})^{\sim}$.

\[
F\overset{\deleted{\unsure{\mathrm{déf}}}}{=}(TC)_{*}. \tag*{\textcircled{\scriptsize $\unsure{e|e'}$}}
\]

\uline{Corollaire} 5.7. Soit \deleted{\unsure{$M$,}} $(W,C,F)$ un triple

\noindent\vrule\hspace{.6em}\begin{minipage}{.93\linewidth}
de Quillen \uline{clos} dans $M$. Pour que
ce triple soit spécial, il f et s.
qu'il satisfasse les conditions suivantes

\textcircled{\scriptsize $a$} $M$ stable par \inserted{petites} $\unsure{\varprojlim|\varinjlim}$,
et ses éléments sont accessibles

\textcircled{\scriptsize $\unsure{c|b},c'$} Il existe une petite partie $C_{0}$ de $C$,
$TC_{0}$ de $TC$, telles que l'on ait
\[
 C\deleted{\unsure{\subset}}={C_{0}}^{\sim},\qquad TC=(TC_{0})^{\sim}.\tag*{$(5.7.1)_{a}$}
\]
\end{minipage}

En effet, toutes les \inserted{autres} conditions \deleted{\unsure{sont}}
a) à e), y inclus c') de 5.6., sont
automatiquement satisfaites.

